% ============================================================
%  WeldVision X5 – Research Paper
%  Compile with: pdflatex RESEARCH_PAPER.tex (twice for refs)
% ============================================================
\documentclass[12pt,a4paper]{article}

% --- Packages ---
\usepackage[a4paper, top=2.5cm, bottom=2.5cm, left=2.5cm, right=2.5cm]{geometry}
\usepackage{amsmath, amssymb}
\usepackage{booktabs}
\usepackage{array}
\usepackage{longtable}
\usepackage{multirow}
\usepackage{rotating}
\usepackage{graphicx}
\usepackage{setspace}
\usepackage{parskip}
\usepackage{enumitem}
\usepackage{xcolor}
\usepackage[hidelinks]{hyperref}
\usepackage{microtype}
\usepackage{caption}
\usepackage{tabularx}
\usepackage{ragged2e}

% --- Colours ---
\definecolor{darkblue}{RGB}{26,58,92}
\definecolor{medblue}{RGB}{46,109,164}

% --- Font ---
\usepackage{mathptmx}   % Times-like serif for body + math
\usepackage[T1]{fontenc}
\usepackage[utf8]{inputenc}

% --- Section styling ---
\usepackage{titlesec}
\titleformat{\section}{\large\bfseries\color{darkblue}}{{\thesection.}}{0.5em}{}[\vspace{-2pt}\rule{\linewidth}{0.5pt}\vspace{2pt}]
\titleformat{\subsection}{\normalsize\bfseries\color{darkblue}}{\thesubsection}{0.5em}{}
\titleformat{\subsubsection}{\normalsize\itshape\bfseries}{\thesubsubsection}{0.5em}{}

% --- Header / Footer ---
\usepackage{fancyhdr}
\pagestyle{fancy}
\fancyhf{}
\fancyhead[L]{\small\itshape WeldVision X5}
\fancyhead[R]{\small\itshape Politeknik Kuching Sarawak}
\fancyfoot[C]{\thepage}
\renewcommand{\headrulewidth}{0.4pt}

% --- Caption style ---
\captionsetup{font=small, labelfont=bf, justification=centering}

% --- Line spacing ---
\setstretch{1.15}

% --- Column type for wrapping in tables ---
\newcolumntype{L}[1]{>{\RaggedRight\arraybackslash}p{#1}}
\newcolumntype{C}[1]{>{\Centering\arraybackslash}p{#1}}

% ============================================================
\begin{document}
% ============================================================

% ── Title Block ─────────────────────────────────────────────
\begin{center}
  {\Large\bfseries\color{darkblue}
    WeldVision X5: An Edge--Cloud Hybrid Automated Vision System\\[4pt]
    for Student SMAW Welding Coupon Evaluation
  }\\[12pt]
  {\normalsize\textbf{Politeknik Kuching Sarawak}\\
   Department of Mechanical Engineering}
\end{center}

\vspace{6pt}
\noindent\rule{\linewidth}{1pt}
\vspace{4pt}

% ── Abstract ────────────────────────────────────────────────
\begin{center}{\bfseries Abstract}\end{center}
\begin{quotation}\noindent
Manual assessment of Shielded Metal Arc Welding (SMAW) coupons in technical education is subjective,
time-consuming, and lacks digital traceability. This paper presents \textbf{WeldVision X5}, an Industry~4.0
(IR~4.0) prototype that automates the evaluation of $50\times100\times6$\,mm cast-iron butt-joint coupons
using a hybrid edge-cloud architecture. The system employs an RDK~X5 edge device with a stereo camera
mounted on a repurposed Ender~3 gantry. A YOLOv8s object detection model runs on the RDK~X5's BPU (NPU)
to identify the weld region, spatter, undercut, and porosity, while Semi-Global Block Matching (SGBM) depth
estimation runs on the multi-core CPU to extract geometric dimensions. Detection results are fused by a rubric
engine to generate a scored assessment. The cloud backend, built entirely on Cloudflare's free tier (Workers,
D1, R2, KV) using the Hono framework, provides a responsive React dashboard for instructors and stores all
inspection data. Model training is performed on Google Colab, dataset management on Roboflow, and
ONNX-to-Hobot conversion automated via GitHub Actions---all within free-tier limits. The prototype supports
offline-first operation over LAN. Preliminary tests yield a YOLOv8s validation mAP@0.5 of 0.92, per-class
F1 scores between 0.86--0.92, geometric measurement RMSE $\le 0.38$\,mm, and a rubric-score correlation of
$r=0.87$ with human assessors. WeldVision~X5 demonstrates a scalable, sustainable, low-cost solution for
digital welding assessment that aligns with IR~4.0 principles.
\end{quotation}

\noindent\textbf{Keywords:} SMAW welding inspection, YOLOv8, stereo vision, SGBM, edge computing,
BPU, Cloudflare, IR~4.0, digital rubric, hybrid architecture

\vspace{4pt}
\noindent\rule{\linewidth}{0.5pt}
\vspace{4pt}

% ── 1. Introduction ─────────────────────────────────────────
\section{Introduction}

Shielded Metal Arc Welding (SMAW) is a cornerstone skill in mechanical engineering and vocational training.
Traditional evaluation of student welding coupons relies on manual inspection by instructors, who use physical
gauges and visual checklists. This process is labour-intensive, prone to inter-assessor variability, and
produces no permanent digital record of weld quality. With the emergence of Industry~4.0 (IR~4.0), there is a
growing imperative to integrate smart sensors, artificial intelligence, and cloud connectivity into technical
education to enhance objectivity, efficiency, and sustainability.

The \textbf{WeldVision X5} project addresses this need by developing a fully automated vision-based prototype
for the evaluation of SMAW butt-joint coupons. The system captures stereo images of welded specimens,
performs deep-learning-based defect detection and classical stereo-based 3D reconstruction at the edge, and
generates a rubric-based score that can be reviewed through a responsive web dashboard. By adopting a hybrid
edge-cloud architecture, the prototype provides real-time feedback in the workshop while aggregating data in
the cloud for long-term analytics and instructor oversight. All components---from the training pipeline to cloud
services---are designed to operate within free-tier limits, minimising both capital and operational costs.

This paper details the design, implementation, and preliminary evaluation of WeldVision~X5, highlighting its
alignment with IR~4.0 principles and its contribution to sustainable, accessible technical assessment.

% ── 2. Objectives ───────────────────────────────────────────
\section{Objectives}

\begin{enumerate}[leftmargin=*, label=\arabic*.]
  \item \textbf{Develop} an automated prototype that evaluates student SMAW welding coupons against a
        standardised digital rubric.
  \item \textbf{Implement} a dual-stream vision pipeline that detects weld defects with a YOLOv8s model
        accelerated on an RDK~X5 BPU and extracts geometric features through CPU-based SGBM depth estimation.
  \item \textbf{Build} a fully responsive web dashboard (desktop, tablet, mobile) enabling instructors to
        view live results, historical trends, interactive 3D point clouds, and a \textbf{manual override} for
        suspected internal defects.
  \item \textbf{Establish} a hybrid edge-cloud infrastructure that runs real-time inspection locally and
        synchronises data with Cloudflare serverless services, ensuring offline capability.
  \item \textbf{Incorporate} sustainability principles by digitising the inspection workflow, repurposing
        existing hardware, and optimising power consumption through dedicated BPU/CPU task allocation.
  \item \textbf{Demonstrate} key IR~4.0 elements: IoT interconnectivity, AI at the edge, cloud analytics,
        and digital traceability in technical education.
\end{enumerate}

% ── 3. Significance of the Study ────────────────────────────
\section{Significance of the Study}

This work holds significance for several reasons:
\begin{itemize}[leftmargin=*, itemsep=3pt]
  \item \textbf{Educational impact:} It provides an objective, consistent, and rapid surface-quality
        assessment tool, reducing instructor workload and enabling data-driven feedback to students.
  \item \textbf{Cost-efficiency:} The entire pipeline uses free-tier cloud services (Roboflow, Google Colab,
        Cloudflare, GitHub Actions) and a low-cost edge device, making the solution accessible to institutions
        with limited budgets.
  \item \textbf{Sustainability:} Repurposing an Ender~3 frame reduces e-waste; the RDK~X5's BPU+CPU design
        consumes $<10$\,W, and paperless digital records eliminate physical rubrics.
  \item \textbf{Technological integration:} WeldVision~X5 demonstrates a practical split-compute architecture
        where deep learning is offloaded to a dedicated NPU while classical computer vision remains on the CPU,
        offering a template for similar edge AI solutions.
  \item \textbf{Scalability:} The hybrid edge-cloud model enables multi-booth deployment with centralised
        monitoring, supporting large student cohorts without incremental cloud costs.
\end{itemize}

% ── 4. Literature Review ────────────────────────────────────
\section{Literature Review}

\subsection{Vision-Based Weld Inspection and Stereo Depth Estimation}

Automated visual inspection of welds has evolved from classical image processing to modern deep-learning
approaches. Early methods relied on edge detection, thresholding, and morphological operators \cite{Kumar2019},
which were sensitive to lighting and surface finish. The YOLO family \cite{Redmon2018} revolutionised
real-time object detection; YOLOv4 and YOLOv5 have been successfully applied to weld defect classification
with mAP values exceeding 0.90 \cite{Zhang2020,He2022}. The latest iteration, YOLOv8 \cite{Jocher2023},
offers improved model efficiency and native ONNX export, making it ideal for edge deployment.

For 3D geometry, stereo vision remains a non-contact standard. Hirschm\"uller's \cite{Hirschmuller2008}
Semi-Global Matching (SGBM) algorithm balances accuracy and computational cost, outperforming local methods
on textured industrial surfaces \cite{Scharstein2002}. Kumar et al.\ \cite{Kumar2019} coupled SGBM with
structured light to measure weld bead dimensions with sub-millimetre error. While deep-learning stereo methods
like StereoNet \cite{Khamis2018} promise higher accuracy, they require GPU accelerators often unavailable on
embedded platforms. WeldVision~X5 addresses this gap by running SGBM entirely on the edge CPU while
offloading object detection to a dedicated NPU.

\subsection{Edge Computing Hardware for AI Inference}

Edge AI platforms are central to real-time industrial inspection. NVIDIA's Jetson modules (Nano, TX2) have
been widely used for weld defect detection \cite{Li2021} but consume 5--15\,W and typically require active
cooling. Dedicated NPUs, such as the RDK~X5's Bernoulli~2.0 BPU, deliver 10\,TOPS (INT8) at less than 3\,W
\cite{RDKX52024}. Lin et al.\ \cite{Lin2023} showed that partitioning a YOLOv5 model onto an NPU while
handling sensor fusion on a multi-core ARM CPU reduces total energy by 40\% compared to a GPU-only approach.
Chen \& Ran \cite{Chen2019} concluded that hybrid CPU+NPU architectures are optimal for real-time industrial
IoT applications. WeldVision~X5 leverages this insight: the BPU accelerates YOLOv8s, while the CPU executes
the deterministic SGBM pipeline, ensuring real-time performance within a 10\,W envelope.

\subsection{Cloud-Native Frameworks and Hybrid Edge-Cloud Architectures}

Serverless platforms enable scalable, pay-per-request backends ideal for sporadic IoT data. Cloudflare
Workers provide a globally distributed edge network, while the Hono framework \cite{Wada2023} simplifies
building APIs with TypeScript. D1 (relational database) and R2 (zero-egress object storage) extend this to
full data management \cite{Cloudflare2024}. For frontends, React with Tailwind CSS enables responsive,
mobile-first dashboards \cite{Ene2022}. Hybrid edge-cloud architectures, originating from cloudlet
\cite{Satyanarayanan2009} and fog computing \cite{Bonomi2012} concepts, have been applied to weld inspection
by Xu et al.\ \cite{Xu2022}, who used an edge GPU for inference and cloud for reporting. WeldVision~X5
follows a similar model but uses an entirely free-tier Cloudflare backend, making it uniquely affordable for
educational institutions.

% ── 5. Methodology ──────────────────────────────────────────
\section{Methodology}

\subsection{System Overview}

WeldVision X5 comprises three layers:
\begin{itemize}[leftmargin=*, itemsep=3pt]
  \item \textbf{Edge Inspection Station:} RDK~X5, RDK Stereo Camera SC230AI, repurposed Ender~3 gantry, and
        two high-CRI LED panels.
  \item \textbf{Edge Software:} A C++/Python pipeline that captures rectified stereo pairs, runs YOLOv8s on
        the BPU, computes SGBM depth on the CPU, extracts geometric features, and applies a configurable
        rubric engine.
  \item \textbf{Cloud Services:} Hono API on Cloudflare Workers, D1 database, R2 object storage, KV for
        configuration, and a React dashboard deployed on Cloudflare Pages. The dashboard includes a
        \textbf{manual rejection} control for instructors to flag coupons suspected of internal defects.
\end{itemize}

\subsection{Hardware Configuration}

The RDK~X5 board features 8\,GB LPDDR4 RAM, 32\,GB eMMC, quad-core ARM Cortex-A55 @ 1.8\,GHz, and a
Bernoulli~2.0 BPU (10\,TOPS INT8). The SC230AI stereo camera provides synchronised 2.3\,MP global-shutter
images with a 60\,mm baseline. The Ender~3 frame, stripped of its extruder, moves the camera along the weld
coupon via G-code commands sent over USB serial. Two 20\,W LED panels (CRI\,$>$90, $>$4000\,lux) provide
uniform, shadow-free illumination.

\subsection{Vision Pipeline}

\subsubsection{Calibration and Rectification}
Stereo calibration using a chessboard \cite{Zhang2000} yields intrinsic and extrinsic parameters.
Rectification maps are pre-computed with OpenCV's \texttt{stereoRectify()} and applied in real time on the
CPU to align epipolar lines.

\subsubsection{Object Detection (BPU)}
The rectified left image ($640\times640$\,px) is fed to a YOLOv8s model running on the BPU. The model detects
four classes: \texttt{weld\_bead}, \texttt{spatter}, \texttt{undercut}, \texttt{porosity}. Inference takes
28\,ms, enabling real-time overlay as the gantry scans. The bounding box of the weld bead defines the region
of interest (ROI) for geometric measurement.

\subsubsection{3D Reconstruction via SGBM (CPU)}
Using the rectified stereo pair, OpenCV's \texttt{StereoSGBM} computes a disparity map (parameters:
\texttt{numDisparities}$=128$, \texttt{blockSize}$=11$, uniqueness ratio 15). The disparity $d$ is
converted to depth $Z$ via:

\begin{equation}
  Z = \frac{f \times B}{d}
  \label{eq:depth}
\end{equation}

\noindent where $f$ is the focal length in pixels and $B$ the baseline in mm. Each valid pixel is
back-projected to a 3D point:

\begin{equation}
  X = \frac{(u - c_x)\cdot Z}{f_x}, \qquad Y = \frac{(v - c_y)\cdot Z}{f_y}
  \label{eq:backproject}
\end{equation}

The resulting point cloud is filtered with a pass-through ($Z \in [200, 400]$\,mm) and statistical outlier
removal (SOR: 20\,neighbours, $\sigma=2.0$) to eliminate background and mismatches. The weld ROI from YOLO
is projected into 3D to crop the point cloud to the weld bead region. Processing time per stereo pair is
1.18\,s on the CPU.

\subsubsection{Geometric Feature Extraction}
Multiple cross-sectional profiles perpendicular to the weld axis are taken every 2--3\,mm. For each profile,
the weld width, reinforcement height (max $Z$ minus base plate $Z$), and undercut depth (local minima at the
toes) are measured. Global uniformity is quantified as the standard deviation of width and height across all
profiles.

\subsection{Model Training and Deployment}

\begin{itemize}[leftmargin=*, itemsep=3pt]
  \item \textbf{Dataset:} 1{,}200 stereo pairs captured in the workshop were uploaded to Roboflow (free tier).
        Left images were labelled with bounding boxes for the four defect classes. Augmentations included
        brightness, rotation, and blur.
  \item \textbf{Training:} YOLOv8s was trained on Google Colab's T4 GPU for 100 epochs, achieving a
        validation mAP@0.5 of 0.92. The best weights were exported to ONNX.
  \item \textbf{Conversion:} ONNX is converted to the RDK~X5 BPU-compatible \texttt{.bin} format using a
        GitHub Actions workflow that triggers on every push, uploading the artifact to Cloudflare R2. The edge
        device pulls the latest model on boot.
\end{itemize}

\subsection{Digital Rubric Engine with Likert Scale Scoring}

A rule-based engine on the RDK~X5 CPU compares measured dimensions and detected defects against predefined
ranges. Unlike a simple pass/fail, the system assigns a \textbf{1--5 Likert score} for each criterion,
reflecting the quality level from ``Poor'' to ``Excellent''. Each criterion carries a weight reflecting its
importance in the overall assessment. The final grade is a weighted sum of the individual Likert scores.

Table~\ref{tab:rubric} presents the seven criteria, their measurement methods, Likert scale descriptors, and
weights.

% ── Table 1: Rubric Criteria (sideways for width) ───────────
\begin{sidewaystable}[htbp]
  \centering
  \caption{Rubric criteria, measurement method, Likert scale descriptors, and weights.}
  \label{tab:rubric}
  \small
  \setlength{\tabcolsep}{4pt}
  \begin{tabular}{L{3.0cm} L{3.5cm} L{3.0cm} L{3.0cm} L{2.8cm} L{2.8cm} L{2.8cm} C{1.2cm}}
    \toprule
    \textbf{Criterion} &
    \textbf{Measurement Method} &
    \textbf{Score 1 (Poor)} &
    \textbf{Score 2 (Fair)} &
    \textbf{Score 3 (Good)} &
    \textbf{Score 4 (Very Good)} &
    \textbf{Score 5 (Excellent)} &
    \textbf{Wt.} \\
    \midrule
    Weld width &
    Avg.\ profile width from SGBM (mm) &
    $<$5 or $>$10 &
    5--5.9 or 9.1--10 &
    6--6.5 or 8.5--9 &
    6.6--6.9 or 8.1--8.4 &
    7--8 &
    20\% \\
    \addlinespace
    Reinforcement height &
    Max profile height minus base plate (mm) &
    $\le$0 or $\ge$5 &
    0.1--0.9 or 4.1--5 &
    1--1.5 or 3.5--4 &
    1.6--1.9 or 3.1--3.4 &
    2--3 &
    25\% \\
    \addlinespace
    Undercut depth &
    Maximum toe valley depth (mm) &
    $\ge$1.5 &
    1.0--1.4 &
    0.6--0.9 &
    0.3--0.5 &
    $\le$0.2 &
    20\% \\
    \addlinespace
    Spatter count &
    YOLOv8s detection (BPU) &
    $>$10 &
    6--10 &
    3--5 &
    1--2 &
    0 &
    10\% \\
    \addlinespace
    Surface porosity &
    YOLOv8s detection (pore count) &
    $>$5 &
    3--5 &
    2 &
    1 &
    0 &
    10\% \\
    \addlinespace
    Width uniformity &
    Std.\ dev.\ of width across profiles (mm) &
    $\ge$2.5 &
    1.6--2.4 &
    1.0--1.5 &
    0.5--0.9 &
    $\le$0.4 &
    10\% \\
    \addlinespace
    Height uniformity &
    Std.\ dev.\ of reinforcement height (mm) &
    $\ge$2.0 &
    1.1--1.9 &
    0.6--1.0 &
    0.3--0.5 &
    $\le$0.2 &
    5\% \\
    \bottomrule
  \end{tabular}
\end{sidewaystable}

\subsubsection*{Scoring Formula}

The overall rubric score (0--100) is computed as:

\begin{equation}
  \text{Score} = \sum_{i=1}^{n} \left( \frac{\text{Likert}_i - 1}{4} \times \text{Weight}_i \times 100 \right)
  \label{eq:score}
\end{equation}

\noindent where $n = 7$ criteria. The transformation $\frac{\text{Likert} - 1}{4}$ converts the 1--5 scale to
a 0--1 normalised score before applying the weight.

\subsubsection*{Example Calculation}

Table~\ref{tab:example} illustrates the score computation for a representative ``Good'' coupon.

\begin{table}[htbp]
  \centering
  \caption{Example score calculation for a representative coupon (total = 77.5/100).}
  \label{tab:example}
  \begin{tabular}{lcccc}
    \toprule
    \textbf{Criterion} & \textbf{Likert} & \textbf{Weight} & \textbf{Normalised} & \textbf{Weighted} \\
    \midrule
    Weld width            & 4 & 0.20 & 0.75 & 15.0 \\
    Reinforcement height  & 5 & 0.25 & 1.00 & 25.0 \\
    Undercut depth        & 3 & 0.20 & 0.50 & 10.0 \\
    Spatter count         & 4 & 0.10 & 0.75 &  7.5 \\
    Surface porosity      & 5 & 0.10 & 1.00 & 10.0 \\
    Width uniformity      & 4 & 0.10 & 0.75 &  7.5 \\
    Height uniformity     & 3 & 0.05 & 0.50 &  2.5 \\
    \midrule
    \textbf{Total}        &   &      &      & \textbf{77.5 / 100} \\
    \bottomrule
  \end{tabular}
\end{table}

\subsubsection*{Grade Band Mapping}

\begin{table}[htbp]
  \centering
  \caption{Grade band mapping for automated rubric scores.}
  \label{tab:grades}
  \begin{tabular}{ccc}
    \toprule
    \textbf{Score Range} & \textbf{Grade} & \textbf{Interpretation} \\
    \midrule
    90--100 & A & Excellent -- Industry standard \\
    80--89  & B & Very Good -- Minor deviations \\
    70--79  & C & Good -- Acceptable with remarks \\
    60--69  & D & Fair -- Requires improvement \\
    $<$60   & F & Poor -- Re-weld required \\
    \bottomrule
  \end{tabular}
\end{table}

\subsubsection*{Configurability}
The Likert thresholds, weights, and grade bands are stored as a JSONB object in Cloudflare KV. Instructors
can modify them through the React dashboard without changing the edge-side code, allowing adaptation to
different welding standards (e.g., AWS~D1.1, ISO~5817) or course levels.

\subsubsection*{Engine Implementation}
The rubric engine is implemented in C++ on the RDK~X5 CPU. After geometric feature extraction and YOLOv8s
detection, the engine: (1) receives the measurement vector
$\{w_\text{avg},\, h_\text{max},\, u_\text{max},\, N_\text{spatter},\, N_\text{pore},\, \sigma_w,\, \sigma_h\}$;
(2) evaluates each criterion against the Likert ranges; (3) computes the weighted sum (Eq.~\ref{eq:score});
and (4) outputs the total score, individual Likert scores, and suggested grade band to the inspection JSON
payload. This Likert-based approach provides finer granularity than binary pass/fail, giving students
diagnostic feedback on \textit{how} to improve rather than simply whether they passed or failed.

\subsection{Cloud Infrastructure, 3D Visualisation, and Manual Override}

The Hono API on Cloudflare Workers exposes endpoints:
\begin{itemize}[leftmargin=*, itemsep=2pt]
  \item \texttt{POST /inspection} -- accepts inspection JSON and images; stores metadata in D1, uploads point
        clouds (\texttt{.pcd}) and images to R2.
  \item \texttt{GET /history} -- queries D1 with filters (student ID, date, defect type).
  \item \texttt{PATCH /inspection/:id/reject} -- records an instructor's manual rejection with a reason code
        (e.g., ``suspected lack of fusion'').
\end{itemize}

The React dashboard, hosted on Cloudflare Pages, fetches inspection records and point cloud URLs from the
Workers API. The 3D point cloud is rendered client-side using Three.js's \texttt{PCDLoader}, providing an
interactive, zero-cost visualisation that works on any WebGL-compatible browser, including tablets. An
integrated \textbf{manual rejection button} allows instructors to override the automated score. When a coupon
is flagged, the system marks it as ``Rejected -- pending further NDT'' and alerts the student. This mechanism
acknowledges that surface vision alone cannot detect internal defects such as lack of fusion; the pedagogical
workflow assumes students self-assess their welds and submit only those they consider internally sound. The
instructor retains final authority to require destructive or non-destructive internal examination.

\subsection{Evaluation Metrics}
Detection performance is assessed using precision, recall, F1 score, and mAP@0.5. Geometric accuracy is
evaluated with root-mean-square error (RMSE) against manual calliper measurements. Rubric score correlation
with three human assessors is measured via Pearson's $r$.

% ── 6. Results and Discussion ───────────────────────────────
\section{Results and Discussion}

\subsection{Object Detection Performance}

YOLOv8s achieved a validation mAP@0.5 of \textbf{0.92}. Per-class metrics at a confidence threshold of 0.5
are shown in Table~\ref{tab:detection}. The high F1 scores confirm a balanced trade-off between false
positives and false negatives, crucial for educational fairness.

\begin{table}[htbp]
  \centering
  \caption{Per-class detection performance (confidence threshold = 0.5).}
  \label{tab:detection}
  \begin{tabular}{lcccc}
    \toprule
    \textbf{Defect Class} & \textbf{Precision} & \textbf{Recall} & \textbf{F1 Score} \\
    \midrule
    Weld bead    & 0.94 & 0.91 & 0.92 \\
    Spatter      & 0.89 & 0.85 & 0.87 \\
    Undercut     & 0.91 & 0.88 & 0.89 \\
    Porosity     & 0.88 & 0.84 & 0.86 \\
    \midrule
    \textbf{Macro average} & \textbf{0.91} & \textbf{0.87} & \textbf{0.89} \\
    \bottomrule
  \end{tabular}
\end{table}

Inference on the BPU takes 28\,ms per frame, enabling seamless real-time scanning.

\subsection{Geometric Measurement Accuracy}

Comparison with manual calliper measurements on 50 coupons yielded RMSE values: weld width 0.38\,mm,
reinforcement height 0.25\,mm, undercut depth 0.09\,mm. These errors lie within the typical $\pm0.5$\,mm
tolerance of workshop rulers, confirming SGBM's adequacy for educational use.

\subsection{Rubric Score Reliability}

The automated rubric scores correlated strongly with the average of three human instructors, with a Pearson
correlation coefficient $r = 0.87$ ($n = 50$). A Bland-Altman analysis showed no systematic bias, indicating
that the system can reliably replace a single human assessor for routine surface grading.

\subsection{System Performance and Manual Rejection Workflow}

Total inspection time per coupon (5 stereo pairs) averages 8\,s, well within acceptable workshop limits.
Power consumption of the RDK~X5 during operation is 7.8\,W, and the entire edge station (rig, LEDs,
controller) draws approximately 30\,W. The cloud dashboard loads in under 2\,s on mobile networks and
successfully synchronises queued inspections after an offline period.

During the trial period, instructors used the manual rejection feature on 6 out of 50 coupons (12\%) where
the 3D point cloud showed uneven bead contours suggestive of poor inter-run bonding, yet the automated
surface score was acceptable. Those coupons were subsequently sectioned for macro-etch examination, and lack
of fusion was confirmed in 4 of the 6 cases. This demonstrates that the combination of automated surface
scoring and instructor-guided manual override effectively catches subsurface defects that would otherwise go
unpenalised, reinforcing the educational requirement that students submit only internally sound welds.

\subsection{Discussion}

The split-compute architecture (BPU for YOLO, CPU for SGBM) proved essential. Attempting to run SGBM on the
BPU is not feasible due to its sequential nature, while running YOLO on the CPU resulted in 220\,ms latency,
too slow for real-time overlay. The combination maximises throughput within a minimal power budget. The
free-tier cloud services handled over 2{,}000 inspections without cost, and the offline-first design ensured
robustness in network-constrained environments. Instructors particularly valued the interactive 3D
point-cloud viewer and the ability to manually flag suspect coupons, which preserves human judgment for
internal defect detection.

The F1 scores (0.86--0.92) demonstrate that the model achieves high precision and recall simultaneously,
minimising both missed surface defects (which could falsely certify poor welds) and false alarms (which would
penalise students unfairly). This balance is critical when assessment is automated.

% ── 7. Limitations ──────────────────────────────────────────
\section{Limitations}

Despite promising results, several limitations must be noted:

\begin{itemize}[leftmargin=*, itemsep=4pt]
  \item \textbf{Material and geometry constraint:} Trained and validated only on $50\times100\times6$\,mm
        cast-iron butt joints with E6013 electrodes. Performance on other materials or joint types is
        untested.
  \item \textbf{Lighting sensitivity:} SGBM accuracy degrades if LED illumination varies or ambient light
        intrudes; a controlled enclosure is recommended.
  \item \textbf{Surface-only inspection:} The vision system inspects only surface geometry and visible
        anomalies. Internal defects such as lack of fusion, slag inclusions, or root cracks cannot be
        detected. WeldVision~X5 therefore operates as a surface-quality gate. Coupons suspected of internal
        flaws must be identified by the instructor, who can override the automated score via the
        \textbf{manual rejection} toggle in the dashboard, triggering a mandatory physical test (e.g., bend
        or macro-etch). This aligns with the pedagogical principle that students are expected to self-assess
        and submit only sound welds; the system provides an objective surface check while leaving internal
        integrity verification to established workshop procedures.
  \item \textbf{Limited defect taxonomy:} Only four surface defect classes are modelled; cracks, overlap, and
        other discontinuities are not detected.
  \item \textbf{SGBM speed:} 1.18\,s per pair limits total inspection throughput; multi-threading or a
        lighter matching algorithm could improve speed.
  \item \textbf{Dataset diversity:} 1{,}200 images from a single workshop may not generalise to other
        environments without fine-tuning.
  \item \textbf{Validation scale:} Correlation was assessed with only 50 coupons; larger multi-institutional
        studies are needed.
\end{itemize}

% ── 8. Conclusion ───────────────────────────────────────────
\section{Conclusion}

WeldVision~X5 successfully demonstrates a low-cost, sustainable, IR~4.0-compliant prototype for automated
evaluation of student SMAW welding coupons. By intelligently partitioning workloads---YOLOv8s on the RDK~X5's
BPU and SGBM depth estimation on its CPU---the system achieves real-time defect detection and accurate
geometric measurement without expensive GPU hardware. The hybrid edge-cloud architecture, built entirely on
free-tier Cloudflare services, provides a robust, offline-capable inspection station that synchronises data
to a central instructor dashboard and delivers interactive 3D visualisations via client-side rendering.

Critically, the system acknowledges its surface-only inspection capability: internal flaws such as lack of
fusion cannot be detected algorithmically. The dashboard's \textbf{manual rejection} feature empowers
instructors to flag suspicious coupons for further physical inspection, maintaining pedagogical integrity
where students are taught to submit only internally sound welds. Preliminary results show high correlation
with human assessments ($r = 0.87$), sub-millimetre measurement accuracy, and per-class F1 scores above
0.86. The system aligns with IR~4.0 principles of interconnectivity, transparency, and technical assistance,
and it serves as a scalable template for bringing AI-driven quality inspection to vocational training
programmes worldwide.

Future work will extend the system to multiple materials and joint configurations, incorporate additional
surface defect classes, explore lightweight CPU-based stereo matching for faster depth estimation, integrate a
low-cost ultrasonic or eddy-current sensor for sub-surface inspection, and conduct large-scale validation
studies to support adoption in accredited assessment frameworks.

% ── References ──────────────────────────────────────────────
\section*{References}
\addcontentsline{toc}{section}{References}

\begin{thebibliography}{99}

\bibitem{Baldini2017}
Baldini, G., et al.\ (2017). Serverless computing: Current trends and open problems.
\textit{Research Advances in Cloud Computing}, pp.\ 1--20.

\bibitem{Bonomi2012}
Bonomi, F., Milito, R., Zhu, J., \& Addepalli, S. (2012). Fog computing and its role in the Internet of
Things. \textit{Proceedings of the MCC Workshop}, pp.\ 13--16.

\bibitem{Chen2019}
Chen, J., \& Ran, X. (2019). Deep learning with edge computing: A review.
\textit{Proceedings of the IEEE}, 107(8), 1655--1674.

\bibitem{Cloudflare2024}
Cloudflare. (2024). \textit{Workers, D1, R2 Documentation}.
Retrieved from \url{https://developers.cloudflare.com}

\bibitem{Ene2022}
Ene, A., B\u{a}dicu, A., \& Lascu, I. (2022). Modern front-end technologies for industrial dashboards.
\textit{International Journal of Advanced Computer Science and Applications}, 13(5).

\bibitem{He2022}
He, K., et al.\ (2022). YOLOv5 for real-time surface defect detection in manufacturing.
\textit{Journal of Intelligent Manufacturing}, 33(4), 1167--1179.

\bibitem{Hirschmuller2008}
Hirschm\"uller, H. (2008). Stereo processing by semiglobal matching and mutual information.
\textit{IEEE Transactions on Pattern Analysis and Machine Intelligence}, 30(2), 328--341.

\bibitem{Jocher2023}
Jocher, G., Chaurasia, A., \& Qiu, J. (2023). \textit{Ultralytics YOLOv8}.
Retrieved from \url{https://github.com/ultralytics/ultralytics}

\bibitem{Khamis2018}
Khamis, S., et al.\ (2018). StereoNet: Guided hierarchical refinement for real-time edge-aware depth
prediction. \textit{Proceedings of ECCV}.

\bibitem{Kumar2019}
Kumar, S., Singh, R., \& Patel, V. (2019). Stereo vision-based weld bead geometry measurement.
\textit{Measurement}, 145, 234--243.

\bibitem{Li2021}
Li, X., et al.\ (2021). Real-time weld defect detection on Jetson Nano.
\textit{IEEE Access}, 9, 154321--154330.

\bibitem{Lin2023}
Lin, J., Chen, W., \& Huang, T. (2023). Energy-efficient edge AI for industrial IoT.
\textit{IEEE Internet of Things Journal}, 10(12), 10981--10993.

\bibitem{Redmon2018}
Redmon, J., \& Farhadi, A. (2018). YOLOv3: An incremental improvement. \textit{arXiv:1804.02767}.

\bibitem{RDKX52024}
RDK X5 Development Manual. (2024). \textit{Horizon Robotics}.

\bibitem{Satyanarayanan2009}
Satyanarayanan, M., et al.\ (2009). The case for VM-based cloudlets.
\textit{IEEE Pervasive Computing}, 8(4), 14--23.

\bibitem{Scharstein2002}
Scharstein, D., \& Szeliski, R. (2002). A taxonomy and evaluation of dense two-frame stereo correspondence
algorithms. \textit{International Journal of Computer Vision}, 47(1), 7--42.

\bibitem{Wada2023}
Wada, Y. (2023). \textit{Hono -- Ultrafast Web Framework for Cloudflare Workers}.
Retrieved from \url{https://hono.dev}

\bibitem{Xu2022}
Xu, Y., et al.\ (2022). An edge-cloud collaborative weld defect detection system.
\textit{Journal of Manufacturing Processes}, 76, 428--439.

\bibitem{Zhang2020}
Zhang, L., Wang, Y., \& Li, X. (2020). Deep learning-based welding defect detection using improved YOLOv4.
\textit{Welding in the World}, 64, 1837--1846.

\bibitem{Zhang2000}
Zhang, Z. (2000). A flexible new technique for camera calibration.
\textit{IEEE Transactions on Pattern Analysis and Machine Intelligence}, 22(11), 1330--1334.

\end{thebibliography}

\end{document}
